\documentclass[11pt,a4paper]{article}
\usepackage{kotex}
\setmainhangulfont{AppleMyungjo}
\setsanshangulfont{Apple SD Gothic Neo}
\setmonofont{Apple SD Gothic Neo}
\usepackage[margin=24mm]{geometry}
\usepackage{longtable,booktabs,array,tabularx}
\usepackage{hyperref}
\usepackage{xurl}
\usepackage{graphicx}
\usepackage{enumitem}
\usepackage{fancyhdr}
\usepackage{setspace}
\usepackage{parskip}
\setlength{\parindent}{0pt}
\setlength{\parskip}{0.45em}
\setlength{\headheight}{14pt}
\setlist[itemize]{leftmargin=1.5em}
\setlist[enumerate]{leftmargin=1.8em}
\newcounter{none}
\providecommand{\tightlist}{%
  \setlength{\itemsep}{0pt}\setlength{\parskip}{0pt}}
\sloppy

\hypersetup{
  colorlinks=true,
  linkcolor=black,
  urlcolor=blue,
  pdftitle={주행 상황 실시간 경고 시스템 별첨 문서집},
  pdfauthor={Mobis PBL Team}
}

\pagestyle{fancy}
\fancyhf{}
\fancyhead[L]{별첨 문서집}
\fancyhead[R]{Driving Alert Project}
\fancyfoot[C]{\thepage}

\begin{document}

\begin{titlepage}
\centering
{\Large 주행 상황 실시간 경고 시스템 \par}
\vspace{0.5em}
{\LARGE 별첨 문서집 \par}
\vspace{0.8em}
{\large Supplementary Appendix Bundle \par}
\vspace{1.6em}
{\normalsize 본 문서집은 소논문 본문을 보강하기 위한 안전, 계약, 검증, 진단 문서를 묶어 정리한 reviewer-facing 별첨 자료이다. \par}
\vspace{1.2em}
\begin{tabular}{ll}
Project & Driving Alert Project \\
Bundle Scope & Governance / Runtime Contracts / Verification / Diagnostic \\
Language & Korean reviewer-facing edition \\
Source Rule & Canonical source paths are preserved in each supplementary note \\
\end{tabular}
\vfill
{\normalsize Workspace: \texttt{driving-alert-workproducts/governance/short-paper/appendix} \par}
\vspace{0.4em}
{\normalsize Generated from local source copies under \texttt{appendix/source/} \par}
\end{titlepage}

\tableofcontents
\clearpage

\section*{S1. 안전 및 추적성 기준}
\addcontentsline{toc}{section}{S1. 안전 및 추적성 기준}
본 별첨은 안전 목표, 위험 사건, ECU naming, CAN ID allocation, Test Matrix를 함께 제시하여 프로젝트의 안전성과 검증 추적성을 보강한다.

\subsection*{원문 기준 경로}
\begin{itemize}
\item \texttt{driving-alert-workproducts/governance/00d\_HARA\_Worksheet.md}
\item \texttt{driving-alert-workproducts/governance/00e\_ECU\_Naming\_Standard.md}
\item \texttt{driving-alert-workproducts/governance/00f\_CAN\_ID\_Allocation\_Standard.md}
\item \texttt{driving-alert-workproducts/governance/00g\_Master\_Test\_Matrix.md}
\end{itemize}

\input{generated/00d_HARA_Worksheet.tex}
\clearpage
\input{generated/00e_ECU_Naming_Standard.tex}
\clearpage
\input{generated/00f_CAN_ID_Allocation_Standard.tex}
\clearpage
\input{generated/00g_Master_Test_Matrix.tex}
\clearpage

\section*{S2. 런타임 계약 세트}
\addcontentsline{toc}{section}{S2. 런타임 계약 세트}
본 별첨은 V2X, ADAS, CGW, 출력 계층의 ownership 경계와 CAN/Ethernet runtime contract를 보강한다.

\subsection*{원문 기준 경로}
\begin{itemize}
\item \texttt{canoe/docs/Kor/contracts/communication-matrix.md}
\item \texttt{canoe/docs/Kor/contracts/owner-route.md}
\item \texttt{canoe/docs/Kor/contracts/layer-separation-policy.md}
\item \texttt{canoe/docs/Kor/contracts/ethernet-interface.md}
\item \texttt{canoe/docs/Kor/contracts/multibus-policy.md}
\item \texttt{canoe/docs/Kor/contracts/panel-sysvar-contract.md}
\item \texttt{canoe/docs/Kor/contracts/ethernet-backbone.md}
\end{itemize}

\input{generated/communication-matrix.tex}
\clearpage
\input{generated/owner-route.tex}
\clearpage
\input{generated/layer-separation-policy.tex}
\clearpage
\input{generated/ethernet-interface.tex}
\clearpage
\input{generated/multibus-policy.tex}
\clearpage
\input{generated/panel-sysvar-contract.tex}
\clearpage
\input{generated/ethernet-backbone.tex}
\clearpage

\section*{S3. 검증 계약 세트}
\addcontentsline{toc}{section}{S3. 검증 계약 세트}
본 별첨은 시험 판정 기준, acceptance criteria, executable asset mapping, execution guide를 제시하여 논문의 검증 구조를 보강한다.

\subsection*{원문 기준 경로}
\begin{itemize}
\item \texttt{canoe/docs/Kor/verification/oracle.md}
\item \texttt{canoe/docs/Kor/verification/acceptance-criteria.md}
\item \texttt{canoe/docs/Kor/verification/test-asset-mapping.md}
\item \texttt{canoe/docs/Kor/verification/execution-guide.md}
\item \texttt{canoe/docs/Kor/verification/evidence-policy.md}
\end{itemize}

\input{generated/oracle.tex}
\clearpage
\input{generated/acceptance-criteria.tex}
\clearpage
\input{generated/test-asset-mapping.tex}
\clearpage
\input{generated/execution-guide.tex}
\clearpage
\input{generated/evidence-policy.tex}
\clearpage

\section*{S4. 진단 및 관측성 세트}
\addcontentsline{toc}{section}{S4. 진단 및 관측성 세트}
본 별첨은 서비스 상태와 경계 조건을 독립적으로 관측하는 진단·관측 구조를 보강한다.

\subsection*{원문 기준 경로}
\begin{itemize}
\item \texttt{canoe/docs/Kor/contracts/diagnostic-matrix.md}
\item \texttt{canoe/docs/Kor/contracts/diagnostic-sysvar-contract.md}
\item \texttt{canoe/docs/Kor/verification/diagnostic-seam-design.md}
\end{itemize}

\input{generated/diagnostic-matrix.tex}
\clearpage
\input{generated/diagnostic-sysvar-contract.tex}
\clearpage
\input{generated/diagnostic-seam-design.tex}

\end{document}
